\subsection{Audit, transparency, and legal grounding}\label{sec:bg-audit}

My headline thread, that an edge controller should produce a provable account of its decisions, sits in a three-way gap in the trust literature. On the attack side, Qu et al.\ \cite{collusion2021} show roughly 6\% of vehicles feeding falsified data to an undefended learning controller: the colluders cut their own waiting time by 92.5\% and raise everyone else's by 62.7\%, and their most dangerous result, that a \emph{calibrated} lie beats a maximal one, motivates my stealth-liar scenario; they name real-time anomaly detection as unbuilt future work. On the integrity side, the blockchain architecture of \cite{nonrepudiation2019} rejects 100\% of spoofed inputs at $\sim$39ms but never names a signature scheme and concedes failure when validators collude. Most fundamentally, input integrity is not decision accountability: it certifies what the controller \emph{heard}, not what it \emph{decided} or why.

My audit design therefore borrows from outside traffic. Certificate Transparency \cite{rfc6962} makes tampering detectable without trusting the log operator, through an append-only Merkle log with inclusion and consistency proofs, at a per-decision cost, a signature and a hash, that fits the edge budget where a validator network does not; my chain inherits the property only once its head is externally witnessed (\S\ref{sec:des-audit}). The legal layer is modest: Dahl et al.\ \cite{dahl2024legal} profile pervasive hallucination when language models answer legal questions, ruling out letting any model assert fault, least of all a 0.6B one, so my system emits a \emph{cited evidence pack} (\S\ref{sec:des-audit}) against a UK-law knowledge base in which authority-side liability is legally weak and arises essentially only from conflicting-green misfeasance, driver-fault doctrines are strong, and emergency-vehicle exemptions are conditional. I am aware of no prior signal-control system that signs and chains individual control decisions, or that mounts a defence-side trust mechanism, though I claim this from a targeted rather than systematic review.
